\chapter{Implémentation}
\label{ch:implementation}

Ce chapitre détaille le code de HOMO-CI~: stack technique, structure
du dépôt, fonctionnement des principales couches.

\section{Stack technique}

\begin{table}[h]
\centering
\caption{Composants techniques de HOMO-CI.}
\label{tab:stack}
\small
\begin{tabularx}{\textwidth}{lX}
\toprule
\textbf{Composant} & \textbf{Choix} \\
\midrule
Langage principal & Python 3.11+ \\
Validation des données & Pydantic v2 (modèles immutables) \\
CLI & Click + Rich (affichage couleur) \\
Stockage local & SQLite (\texttt{sqlite3} stdlib) + JSON Lines \\
SDK Azure & \texttt{azure-identity}, \texttt{azure-mgmt-network}, \texttt{azure-mgmt-resource} \\
Scan CVE & Trivy (binaire externe, appel \texttt{subprocess}) \\
Templating & Jinja2 \\
Conversion PDF & Pandoc + WeasyPrint (CSS) \\
Stockage distant & Azure Blob Storage (publication) \\
Orchestration locale & Makefile + Bash \\
CI/CD & GitHub Actions (workflows quotidiens) \\
\bottomrule
\end{tabularx}
\end{table}

\section{Structure du dépôt}

\begin{lstlisting}[caption={Arborescence du paquet
\texttt{homo\_ci}.}]
pipeline/
+-- homo_ci/                  Package Python
|   +-- models.py             Evidence, Change, DossierSnapshot
|   +-- storage.py            SQLite store + JSONL trail
|   +-- watchers/             Collecteurs
|   |   +-- base.py           Interface BaseWatcher
|   |   +-- wc_nsg.py         Azure NSG
|   |   +-- wc_cve.py         Trivy / CVE Docker
|   |   +-- wc_ad.py          Active Directory
|   |   +-- wc_siem.py        Wazuh
|   +-- diff.py               Moteur de diff
|   +-- render.py             Renderer Jinja2 -> Markdown
|   +-- sign.py               Hash chain + signature
|   +-- cli.py                Point d'entree Click
|   +-- tests/                Tests unitaires
+-- templates/                dossier.md.j2
+-- styles/                   CSS pour WeasyPrint
+-- scripts/                  Scripts d'orchestration
+-- Makefile                  Cibles : run, test, clean
+-- pyproject.toml
+-- requirements.txt
\end{lstlisting}

\section{Les watchers}

Tous les watchers héritent de \texttt{BaseWatcher} qui définit un
contrat simple~:

\begin{lstlisting}[language=Python,caption={Interface
\texttt{BaseWatcher}.}]
class BaseWatcher(ABC):
    source: SourceType  # défini par la classe fille

    def __init__(self, store: EvidenceStore) -> None:
        self.store = store
        self.log = logging.getLogger(self.__class__.__name__)

    @abstractmethod
    def collect(self) -> list[Evidence]: ...

    def run(self) -> WatcherResult:
        """Collecte, persiste, retourne un résumé."""
        ...
\end{lstlisting}

Deux watchers sont implémentés en v0.1 (\texttt{wc\_nsg} et
\texttt{wc\_cve}), les autres (\texttt{wc\_ad}, \texttt{wc\_siem})
sont câblés mais en stub.

\subsection{Watcher NSG : exemple complet}

\texttt{wc\_nsg.py} interroge Azure Resource Manager via le SDK
officiel et produit une \texttt{Evidence} par règle NSG~:

\begin{lstlisting}[language=Python,caption={Cœur du watcher NSG.}]
def collect(self) -> list[Evidence]:
    cred = DefaultAzureCredential(
        exclude_interactive_browser_credential=True
    )
    client = NetworkManagementClient(cred, self.subscription_id)
    evidences: list[Evidence] = []
    for nsg in client.network_security_groups.list(
        self.resource_group
    ):
        for rule in nsg.security_rules or []:
            payload = self._serialize_rule(nsg.name, rule)
            ref = f"{nsg.name}/{rule.name}"
            materiality = self._classify_materiality(rule)
            evidences.append(
                Evidence.from_payload(
                    source=self.source,
                    category=Category.NSG,
                    ref=ref,
                    payload=payload,
                    scope=self.SCOPE,
                    weight=Weight.AUTO,
                    materiality=materiality,
                )
            )
    return evidences
\end{lstlisting}

Deux points à noter~:

\begin{itemize}
  \item L'\texttt{Evidence} est construite via
    \texttt{Evidence.from\_payload(...)} qui calcule
    automatiquement le hash SHA-256 sur le payload JSON
    canonique (clés triées, séparateurs compacts). C'est ce hash
    qui permet la déduplication et la détection de changement.
  \item La matérialité est classifiée en HIGH si la règle est
    \texttt{Allow} et autorise du trafic depuis Internet. C'est
    une heuristique simple mais déjà utile pour prioriser les
    alertes.
\end{itemize}

\subsection{Watcher CVE}

\texttt{wc\_cve.py} appelle Trivy en sous-processus~:

\begin{lstlisting}[language=bash]
trivy image --quiet --format json \
    --severity CRITICAL,HIGH,MEDIUM,LOW \
    internal-apps/maformation:latest
\end{lstlisting}

Le JSON de sortie est parsé, et chaque CVE devient une preuve avec
sa sévérité mappée vers \texttt{Materiality}~:

\begin{lstlisting}[language=Python]
_SEV_MAP = {
    "CRITICAL": Materiality.HIGH,
    "HIGH":     Materiality.HIGH,
    "MEDIUM":   Materiality.MEDIUM,
    "LOW":      Materiality.LOW,
}
\end{lstlisting}

Si Trivy n'est pas disponible (\texttt{shutil.which} échoue), le
watcher émet un warning et retourne une liste vide --- pas
d'exception, le pipeline continue.

\section{Le stockage}

Deux supports complémentaires~:

\begin{description}
  \item[SQLite.] Une seule table \texttt{evidence} avec les
    colonnes (\texttt{id, ts, source, category, ref,
    payload\_hash, payload\_json, materiality, weight}). Indexée
    sur \texttt{(source, ref)}. Permet les requêtes rapides du
    diff et du rendu.
  \item[JSONL.] À chaque collecte, un fichier
    \texttt{snapshot-AAAAMMJJ-HHMMSS.jsonl} est écrit avec une
    preuve par ligne. C'est le format d'archive long terme et
    de comparaison entre exécutions.
\end{description}

\section{Le moteur de diff}

Algorithme en $O(N+M)$~: on construit deux dictionnaires indexés
par \texttt{(source, ref)} pour le snapshot précédent et le
courant, puis on compare~:

\begin{lstlisting}[language=Python,caption={Logique de diff.}]
def diff(previous: dict, current: dict) -> list[Change]:
    changes = []
    prev_keys = set(previous.keys())
    curr_keys = set(current.keys())
    for key in curr_keys - prev_keys:
        changes.append(Change(type=ChangeType.CREATED, ...))
    for key in prev_keys - curr_keys:
        changes.append(Change(type=ChangeType.DELETED, ...))
    for key in prev_keys & curr_keys:
        if previous[key].payload_hash != current[key].payload_hash:
            changes.append(Change(type=ChangeType.UPDATED, ...))
    return changes
\end{lstlisting}

Comparaison par hash plutôt que par contenu~: rapide, et
correcte tant que la canonicalisation JSON reste stable.

\section{Le rendu du dossier}

Un seul template Jinja2 (\texttt{templates/dossier.md.j2}), qui
itère sur les preuves groupées par section ANSSI~:

\begin{lstlisting}[language={[LaTeX]TeX},caption={Extrait du
template.}]
# Dossier d'homologation -- {{ snapshot.version }}

Généré le {{ snapshot.generated_at }}
Hash : {{ snapshot.pdf_hash }}

## Étape 4 -- Cartographie du SI

### Règles réseau (NSG)
{% for ev in evidences | selectattr("category.value", "==", "nsg") %}
- {{ ev.ref }} : {{ ev.payload.action }} ...
{% endfor %}
\end{lstlisting}

Le Markdown produit est ensuite converti en PDF par pandoc, avec
une feuille de style CSS pour la mise en page (logo, marges,
police).

\section{La signature et la \emph{hash chain}}

Le hash SHA-256 du PDF généré est calculé après écriture. Chaque
\texttt{DossierSnapshot} référence par \texttt{parent\_hash} le
hash du dossier précédent~: c'est ce qui constitue la \emph{hash
chain}, et qui permet de détecter une falsification rétroactive
d'un dossier antérieur.

\section{Le CLI}

Quatre commandes~:

\begin{lstlisting}[language=bash]
# Collecter les preuves (par défaut nsg + cve)
homo-ci collect --include nsg --include cve

# Comparer avec le snapshot précédent
homo-ci diff

# Générer le dossier (Markdown + PDF)
homo-ci render

# Publier sur Azure Blob (optionnel)
homo-ci publish
\end{lstlisting}

L'enchaînement complet est encapsulé dans une cible Makefile~:

\begin{lstlisting}[language=bash]
make pipeline-run
# Résultat : build/dossier-AAAAMMJJ-HHMMSS.pdf
\end{lstlisting}

\section{Tests}

Tests unitaires \texttt{pytest} dans \texttt{homo\_ci/tests/}, qui
couvrent~:

\begin{itemize}
  \item La canonicalisation et le hash des payloads (déterminisme).
  \item Le diff entre deux ensembles connus.
  \item La sérialisation / désérialisation JSONL.
  \item Le rendu Jinja2 d'un mini-dossier de référence.
\end{itemize}

Le watcher NSG est testé avec un \emph{mock} du SDK Azure (pas de
vrai appel pendant les tests). Le watcher CVE est testé avec un
JSON Trivy figé.

\begin{takeaway}
HOMO-CI est un pipeline Python de taille modeste (environ
\textasciitilde 2~000 lignes hors tests). Le choix des briques
(Pydantic, SQLite, Jinja2, Trivy, SDK Azure officiel) privilégie
la stabilité et la lisibilité à toute prouesse technique.
\end{takeaway}
